% === Begin Label Data ===
% ID: 885
% 难度星级: 
% 标签: 
% 备注: 
% 组卷引用次数: 0
% === End  Label Data ===

\begin{problem}{2024}{G}{上海卷}{21}{导数}
设 $D$ 为 $\mathbb{R}$ 的非空集合，$a,b \in \mathbb{R}$，对于定义域为 $D$ 的函数 $y=f(x)$ 以及点 $M(a,b)$，定义 $s(x)=(x-a)^2+(f(x)-b)^2$，若存在 $x_0 \in D$ 使得 $s(x_0)$ 是函数 $y=s(x)$ 的最小值，则称点 $P(x_0, f(x_0))$ 为函数 $y=f(x)$ 的“$M(a,b)$ 最近点”.

（1）设 $f(x)=\dfrac{1}{x} (x>0)$，$M(0,0)$. 证明：函数 $y=f(x)$ 存在 $M(0,0)$ “最近点”；

（2）设 $f(x)=e^x$，$M(1,0)$，若曲线 $y=f(x)$ 在点 $P(x_0, f(x_0))$ 的切线与直线 $PM$ 垂直，证明 $P(x_0, f(x_0))$ 是函数 $y=f(x)$ 的 $M(1,0)$ “最近点”；

（3）设定域为 $\mathbb{R}$ 的函数 $y=f(x)$ 存在导函数 $y=f'(x)$，且定义域为 $\mathbb{R}$ 的函数 $y=g(x)$ 函数值恒为正数. 若对任意实数 $t \in \mathbb{R}$，总存在曲线 $y=f(x)$ 上的点 $P$，使得点 $P$ 既是函数 $y=f(x)$ 的 $M_1(t-1, f(t)-g(t))$ 最近点，也是函数 $y=f(x)$ 的 $M_2(t+1, f(t)+g(t))$ 最近点，判断 $y=f(x)$ 的单调性，说明理由.
\end{problem}

\begin{answer}
\end{answer}

\begin{solutions}
\end{solutions}